\subsection{Markov decomposition and inverse-map sector factorization}

For the simple MPDO case in Appendix~C.2 of \cite{Cirac2016MPDO_arXiv},
the strong-area-law hypothesis is used locally through the equality case of
strong subadditivity for a normalized three-site reduced state.

\begin{definition}[Local $\eta$-structure for a three-site reduced state]
    \label{def:mpdo_eta_structure}
    \lean{MPOTensor.EtaStructure}
    \leanok
    \uses{def:entropy_quantum_markov_decomposition}
    For a three-site density operator $\rho_{ABC}$, the local $\eta$-structure
    is the quantum Markov decomposition on the middle subsystem $B$ produced by
    equality in strong subadditivity. Concretely, after a unitary change of
    basis on $B$, one has a finite direct-sum decomposition
    \begin{align}
        B
        &=\bigoplus_k(b_1^{(k)}\otimes b_2^{(k)}),
        \notag\\
        \rho_{ABC}
        &=\bigoplus_k
          \rho_{A b_1}^{(k)}\otimes\rho_{b_2 C}^{(k)}.
        \notag
    \end{align}
\end{definition}

\begin{definition}[Canonical projections of a Hayashi decomposition]
    \label{def:mpdo_hayashi_sector_projection}
    \lean{MPOTensor.EtaStructure.coordinateSectorProjection}
    \lean{MPOTensor.EtaStructure.sectorProjection}
    \leanok
    \uses{def:mpdo_eta_structure}
    Let $B=\bigoplus_k(b_1^{(k)}\otimes b_2^{(k)})$ be the
    middle-site decomposition, and let $E_k$ be the coordinate projection onto
    its $k$-th summand. In the original physical basis, define
    $Q_k=U_B^*E_kU_B$. These are the projectors introduced in
    \cite[Appendix~C.2, lines~1364--1368]{Cirac2016MPDO_arXiv}; the unitary
    implements the basis change described in
    \cite[Appendix~C.2, lines~1437--1440]{Cirac2016MPDO_arXiv}.
\end{definition}

\begin{theorem}[Resolution of the identity by Hayashi-sector projections]
    \label{thm:mpdo_hayashi_sector_projection_resolution}
    \lean{MPOTensor.EtaStructure.sectorProjection_isOrthogonal}
    \lean{MPOTensor.EtaStructure.sum_sectorProjection}
    \lean{MPOTensor.EtaStructure.sectorProjection_mul_eq_zero}
    \leanok
    \uses{def:orthogonal_projection_channel,
        def:mpdo_hayashi_sector_projection}
    Each $Q_k$ is an orthogonal projection, distinct summands are mutually
    orthogonal, and the family resolves the identity:
    \begin{align}
        Q_kQ_\ell&=0 \quad\text{if } k\ne\ell,
        \notag\\
        \sum_k Q_k&=\Id.
        \notag
    \end{align}
    The first two assertions follow from the orthogonal direct-sum
    decomposition in
    \cite[Appendix~C.2, lines~1364--1368]{Cirac2016MPDO_arXiv}; the resolution
    of the identity is stated in
    \cite[Appendix~C.2, lines~1693--1697]{Cirac2016MPDO_arXiv}.
\end{theorem}

\begin{proof}
    \leanok
    \uses{def:mpdo_hayashi_sector_projection,
        lem:orthogonal_projection_mul_eq_zero_of_sum_eq_one,
        thm:star_projection_coisometric_transport,
        thm:orthogonal_projection_star_projection_channel}
    The coordinate projections satisfy
    \begin{align}
        E_k^*&=E_k,
        \notag\\
        E_k^2&=E_k,
        \notag\\
        \sum_k E_k&=\Id.
        \notag
    \end{align}
    Reindexing the direct sum and conjugating by $U_B$ preserve these
    identities. Hence each $Q_k=U_B^*E_kU_B$ is a star projection and
    therefore an orthogonal projection, while
    $\sum_k Q_k=U_B^*(\sum_kE_k)U_B=\Id$. Finally, multiplying the resolution
    on both sides by $Q_k$ gives
    \begin{align}
        \sum_\ell Q_kQ_\ell Q_k&=Q_k,
        \notag\\
        Q_kQ_\ell Q_k&=(Q_\ell Q_k)^*(Q_\ell Q_k)\geq0.
        \notag
    \end{align}
    The $\ell=k$ term already equals $Q_k$, so every term with $\ell\ne k$
    vanishes. Thus $(Q_\ell Q_k)^*(Q_\ell Q_k)=0$, hence
    $Q_\ell Q_k=0$. Taking adjoints gives $Q_kQ_\ell=0$ for $\ell\ne k$.
\end{proof}

\begin{definition}[Markov factors of a simultaneous inverse]
    \label{def:mpdo_bnt_markov_factors}
    \lean{MPOTensor.bntMarkovLeftFactor}
    \lean{MPOTensor.bntMarkovRightFactor}
    \leanok
    \uses{def:mpdo_three_site_family_closure, def:mpdo_eta_structure}
    Let $C_x$ and $C_y$ be two rows of a simultaneous left inverse, and
    choose a Hayashi decomposition of the three-site operator. In its
    $k$-th Markov sector, set
    \begin{align}
        A_x^{(k)}(l,l')
        &=\sum_{i_1,j_1} C_{x,(i_1,j_1)}
          \rho_{A b_1;(i_1,l),(j_1,l')}^{(k)},
        \notag\\
        B_y^{(k)}(r,r')
        &=\sum_{i_3,j_3} C_{y,(i_3,j_3)}
          \rho_{b_2 C;(r,i_3),(r',j_3)}^{(k)}.
        \notag
    \end{align}
    These are the two factors obtained in the projected contraction
    of~\cite[Appendix~C.2, lines~1719--1725]{Cirac2016MPDO_arXiv}.
\end{definition}

\begin{theorem}[Outer inverse contraction in Markov coordinates]
    \label{thm:mpdo_bnt_outer_inverse_markov_block}
    \lean{MPOTensor.outerInverseContraction_markov_block}
    \leanok
    \uses{def:mpdo_bnt_markov_factors}
    Let $\Gamma_{x,y}$ denote the outer inverse contraction. In the basis of
    the chosen Hayashi decomposition,
    \begin{align}
        {[U_B\Gamma_{x,y}U_B^\dagger]}_
          {(k,l,r),(k',l',r')}
        &=
        \begin{cases}
            p_k A_x^{(k)}(l,l')B_y^{(k)}(r,r'), & k=k',\\
            0, & k\ne k'.
        \end{cases}
        \notag
    \end{align}
    In particular, the probability $p_k$ remains explicit. This identity
    requires neither a nonvanishing condition on $p_k$ nor assumptions of
    saturated area law or zero correlation length.
\end{theorem}

\begin{proof}\leanok
    Expand the two matrix products and the outer inverse contraction. The
    Hayashi block decomposition makes the middle-site contraction zero for
    $k\ne k'$. For $k=k'$, the remaining four outer sums separate into the
    product defining $A_x^{(k)}$ and $B_y^{(k)}$, with the common factor $p_k$.
\end{proof}

\begin{theorem}[Entrywise BNT--Markov identity]
    \label{thm:mpdo_bnt_markov_key_formula}
    \lean{MPOTensor.isMPOBlockLeftInverse_bnt_markov_key_formula}
    \leanok
    \uses{thm:mpdo_bnt_three_site_outer_inverse_collapse,
      thm:mpdo_bnt_outer_inverse_markov_block}
    Let $x=(s,\alpha_1,\beta_1)$ and
    $y=(t,\alpha_3,\beta_3)$. Comparing the two expressions for the same
    outer inverse contraction gives
    \begin{align}
        &
        \begin{cases}
            p_k A_x^{(k)}(l,l')B_y^{(k)}(r,r'), & k=k',\\
            0, & k\ne k'
        \end{cases}
        \notag\\
        &\quad=
        \begin{cases}
            {(R_s)}_{\beta_3,\alpha_1}
            [U_B\kappa^{(s)}_{\beta_1,\alpha_3}U_B^\dagger]_
              {(k,l,r),(k',l',r')}, & s=t,\\
            0, & s\ne t.
        \end{cases}
        \notag
    \end{align}
    Here
    $\kappa^{(s)}_{\beta_1,\alpha_3}(i,j)
      ={(\mathcal K_s^{ij})}_{\beta_1,\alpha_3}$.
    Thus both the normal-sector and Markov-sector off-diagonal entries are
    retained in the statement, rather than being suppressed by a prior
    choice of equal labels. This is the full entrywise identity used
    in~\cite[Appendix~C.2, lines~1719--1732]{Cirac2016MPDO_arXiv}.

    The closing entry ${(R_s)}_{\beta_3,\alpha_1}$ is the corrected tail-index
    orientation of the display at lines~1422--1438, recorded in
    \path{docs/paper-gaps/cpgsv17_mpdo_sal_zcl_eta_local_structure.tex}.
\end{theorem}

\begin{proof}\leanok
    The preceding theorem identifies the left side with the transformed outer
    inverse contraction. The collapse theorem identifies the contraction
    with zero when $s\ne t$ and with
    ${(R_s)}_{\beta_3,\alpha_1}
      \kappa^{(s)}_{\beta_1,\alpha_3}$ when $s=t$.
    Conjugating this matrix by $U_B$ gives the right side.
\end{proof}

\begin{theorem}[Off-diagonal Markov blocks of a BNT sector]
    \label{thm:mpdo_bnt_markov_offdiagonal}
    \lean{MPOTensor.isMPOBlockLeftInverse_bnt_markov_offdiagonal}
    \leanok
    \uses{thm:mpdo_bnt_markov_key_formula}
    Fix a normal sector $s$ and virtual indices $\alpha_1,\beta_3$ for which
    ${(R_s)}_{\beta_3,\alpha_1}\ne0$. Then, for every $\beta_1,\alpha_3$
    and every pair of distinct Markov sectors $k\ne k'$, one has
    \begin{align}
        [U_B\kappa^{(s)}_{\beta_1,\alpha_3}U_B^\dagger]_
          {(k,l,r),(k',l',r')}
        &=0.
        \notag
    \end{align}
    Thus each physical slice of the $s$-th normal tensor is block diagonal in
    the Markov decomposition. This is equation~\textup{Qks}
    in~\cite[Appendix~C.2, lines~1682--1688 and 1730--1735]
      {Cirac2016MPDO_arXiv}.

    The nonzero closing entry uses the corrected orientation
    ${(R_s)}_{\beta_3,\alpha_1}$ recorded in
    \path{docs/paper-gaps/cpgsv17_mpdo_sal_zcl_eta_local_structure.tex}.
\end{theorem}

\begin{proof}\leanok
    In the entrywise BNT--Markov identity, take equal normal-sector labels and
    distinct Markov-sector labels. The left side vanishes, while the right
    side is the selected closing entry times the displayed block. Cancel the
    nonzero closing entry.
\end{proof}

\begin{definition}[Diagonal BNT block in a Markov sector]
    \label{def:mpdo_bnt_markov_block}
    \lean{MPOTensor.bntMarkovBlock}
    \lean{MPOTensor.BNTMarkovBlockNonzero}
    \leanok
    For a normal sector $s$ and a Markov sector $k$, let
    \begin{align}
        O_s^{(k)}(\beta,\alpha)
        &=
        [U_B\kappa^{(s)}_{\beta,\alpha}U_B^\dagger]_{k,k}.
        \notag
    \end{align}
    The assertion $O_s^{(k)}\ne0$ means that this block is nonzero for at
    least one pair of virtual indices $\beta,\alpha$.
\end{definition}

\begin{theorem}[Existence of a normal-sector label]
    \label{thm:mpdo_bnt_markov_label_exists}
    \lean{MPOTensor.exists_bntMarkovBlockNonzero_of_probability_ne_zero}
    \leanok
    \uses{def:mpdo_bnt_markov_block}
    If $p_k\ne0$, then there is a normal sector $s$ for which
    $O_s^{(k)}\ne0$.

    This is the existence part of equation~\textup{QkKjs} in
    \cite[Appendix~C.2, lines~1733--1737]{Cirac2016MPDO_arXiv}. Restricting
    to $p_k\ne0$ is the literal version of the source's preceding removal of
    summands for which $Q_k\sigma_3^{(N)}Q_k=0$.
\end{theorem}

\begin{proof}\leanok
    Choose nonzero diagonal entries of the trace-one left and right states in
    sector $k$. The corresponding diagonal entry of the Hayashi
    decomposition is nonzero because $p_k\ne0$. Expanding the same entry by
    the three-site normal-sector family therefore leaves at least one nonzero
    summand, whose middle tensor has $O_s^{(k)}\ne0$.
\end{proof}

\begin{theorem}[Uniqueness from nonzero closing matrices]
    \label{thm:mpdo_bnt_markov_label_unique_nonzero_closing}
    \lean{MPOTensor.bntMarkovBlockNonzero_eq_of_probability_ne_zero}
    \lean{MPOTensor.existsUnique_bntMarkovBlockNonzero_of_probability_ne_zero}
    \leanok
    \uses{thm:mpdo_bnt_markov_key_formula,
      thm:mpdo_bnt_markov_label_exists,def:mpdo_bnt_markov_block}
    Suppose, in addition, that every closing matrix $R_s$ is nonzero. If
    $p_k\ne0$, there is a unique normal sector $j_k$ for which
    $O_{j_k}^{(k)}\ne0$.

    This is the cancellation argument in equation~\textup{QkKjs}. The
    stated nonzero-closing-matrix condition isolates the step obtained in the
    source from simplicity and the common-weight choice at
    \cite[Appendix~C.2, lines~1714--1718]{Cirac2016MPDO_arXiv}. Its derivation
    from those preceding hypotheses is
    Theorem~\ref{thm:mpdo_normalized_three_site_bnt_family_closure}.
\end{theorem}

\begin{proof}\leanok
    Assume that sectors $s\ne t$ both have nonzero $k$-blocks. A diagonal
    instance of the BNT--Markov identity for $s$ gives a nonzero left factor.
    A mixed instance with normal-sector labels $s,t$ forces the right factor
    selected from $t$ to vanish. A diagonal instance for $t$ then has zero
    left-hand side and a nonzero right-hand side, a contradiction.
\end{proof}

\begin{definition}[Normal-sector projections on the active support]
    \label{def:mpdo_bnt_sector_projection}
    \lean{MPOTensor.bntSectorLabel}
    \lean{MPOTensor.bntSectorProjection}
    \lean{MPOTensor.activeMarkovProjection}
    \leanok
    Under the hypotheses of
    Theorem~\ref{thm:mpdo_bnt_markov_label_unique_nonzero_closing}, define
    \begin{align}
        P_s
        &=\sum_{\substack{k:p_k\ne0\\j_k=s}}Q_k,
        \notag\\
        P_{\mathrm{act}}
        &=\sum_{k:p_k\ne0}Q_k.
        \notag
    \end{align}
    This is equation~\textup{Pis}, restricted to the positive-weight Markov
    support retained here.
\end{definition}

\begin{theorem}[Orthogonality and resolution on the active support]
    \label{thm:mpdo_bnt_sector_projection_resolution}
    \lean{MPOTensor.bntSectorProjection_isOrthogonal}
    \lean{MPOTensor.bntSectorProjection_mul_eq_zero}
    \lean{MPOTensor.sum_bntSectorProjection}
    \leanok
    \uses{def:mpdo_bnt_sector_projection}
    Every $P_s$ is an orthogonal projection, distinct $P_s$ are mutually
    orthogonal, and $\sum_s P_s=P_{\mathrm{act}}$.
\end{theorem}

\begin{proof}\leanok
    Each $P_s$ is a sum of mutually orthogonal Markov-sector projections.
    Distinct labels have disjoint sums. Finally, summing over $s$ counts each
    positive-weight Markov sector exactly once.
\end{proof}

\begin{theorem}[SAL and arbitrary-cut equality in strong subadditivity]
    \label{thm:mpdo_sal_ssa_equality}
    \lean{MPOTensor.isSSAEquality_tripartite_m_of_isSAL}
    \lean{MPOTensor.isSSAEquality_tripartite_of_isSAL}
    \lean{MPOTensor.isSSAEquality_threeSite_of_isSAL}
    \lean{MPOTensor.%
      mutualInfoChain_one_eq_of_isSSAEquality_tripartite_m}
    \lean{MPOTensor.isSAL_of_isSSAEquality_tripartite_m}
    \leanok
    \uses{def:is_sal, def:ssa_equality}
    Let $1\leq m\leq\lfloor N/2\rfloor$, and divide the chain into four
    consecutive regions $A,B,C,D$ of lengths $m-1,1,N-m-1,1$. If the tensor
    saturates the area law, then its marginal on $ABC$ satisfies
    $S(ABC)+S(B)=S(AB)+S(BC)$. Indeed, this is the entropy form of
    $I_1=I_m$. Taking $m=2$ gives the regions of lengths $1,1,N-3$, and then
    taking $N=4$ shows that the three-site reduced state of the four-site
    chain satisfies equality in strong subadditivity.

    Conversely, suppose the tensor generates positive semidefinite periodic
    operators and their traces are nonzero at every positive chain length.
    If the displayed equality holds for every $N$ and every
    $1\leq m\leq\lfloor N/2\rfloor$, then the tensor saturates the area law.
    This is the choice of regions used in
    \cite[Appendix~C.2, lines~1760--1780]{Cirac2016MPDO_arXiv}.
    The endpoint corrects the strict inequality printed at source line~1771:
    Definition~4.6 and the concluding comparison both require
    $m=\lfloor N/2\rfloor$. This local correction is recorded in
    \path{docs/paper-gaps/cpgsv17_mpdo_sal_zcl_eta_local_structure.tex}.
\end{theorem}

\begin{proof}
    \leanok
    \uses{thm:mpdo_tripartite_block_entropy_identifications}
    By saturation, $I_1=I_m$. Expanding the two mutual informations and
    cancelling the entropy of the full chain gives
    $S_{N-1}+S_1=S_m+S_{N-m}$, which is the stated equality for the regions
    $A,B,C$. Conversely, this equality gives $I_1=I_m$. If it holds for every
    admissible $m$, comparing the cases $m=L$ and $m=L+1$ gives
    $I_L=I_{L+1}$ and hence saturation of the area law.
\end{proof}

\begin{theorem}[All-cut quantum Markov structure implies SAL]
    \label{thm:mpdo_all_cut_markov_implies_sal}
    \lean{MPOTensor.isSAL_of_quantumMarkovDecomposition_tripartite_m}
    \leanok
    \uses{def:is_sal, def:mpdo, def:entropy_quantum_markov_decomposition}
    Suppose that the periodic operators are positive semidefinite and have
    nonzero trace at every positive chain length. For every $N$ and every
    $1\leq m\leq\lfloor N/2\rfloor$, divide the chain into four consecutive
    regions of lengths $m-1,1,N-m-1,1$. If the marginal on the first three
    regions admits a quantum Markov decomposition on its one-site middle
    system, then the tensor saturates the area law.
    % This is the final entropy step at source lines 1606--1619. The missing
    % implication from its commuting-bond and source-ZCL hypotheses is recorded in
    % docs/paper-gaps/cpsv16_commuting_form_to_sal.tex.
\end{theorem}

\begin{proof}\leanok
    \uses{thm:mpdo_sal_ssa_equality,
        thm:entropy_ssa_equality_quantum_markov}
    The quantum Markov decomposition gives equality in strong subadditivity
    at every admissible cut:
    $S(\rho_{ABC})+S(\rho_B)=S(\rho_{AB})+S(\rho_{BC})$.
    These equalities at every admissible cut imply saturation of the area law
    by Theorem~\ref{thm:mpdo_sal_ssa_equality}.
\end{proof}

\begin{theorem}[SAL implies local $\eta$-structure]
    \label{thm:sal_implies_eta_structure}
    \lean{MPOTensor.sal_implies_eta_structure}
    \leanok
    \uses{def:ssa_equality, def:mpdo_eta_structure}
    If a normalized three-site reduced state satisfies the equality case of
    strong subadditivity --- the local entropy form of the simple-MPDO strong
    area law --- then it admits the local $\eta$-structure of
    Definition~\ref{def:mpdo_eta_structure}.
\end{theorem}

\begin{proof}
    \leanok
    \uses{thm:entropy_ssa_equality_quantum_markov}
    This is the forward implication of the Hayashi equality
    characterization. The product-reference raw Petz formula, including its
    singular-support compression, is given by
    Theorem~\ref{thm:entropy_petz_product_tensor}. The singular-support
    equality-to-recovery theorem is still separate. Beyond it, the remaining
    structural ingredient is the family-level Koashi--Imoto decomposition and
    the induced action of the middle-system channel on its common direct-sum
    factors.
\end{proof}

\begin{lemma}[Partial trace along a right tensor factor]
    \label{lem:partial_trace_right_along}
    \lean{Matrix.partialTraceRightAlong,
      Matrix.partialTraceRightAlong_apply,
      Matrix.PosSemidef.partialTraceRightAlong,
      Matrix.trace_partialTraceRightAlong}
    \leanok
    Suppose that $C\simeq C'\otimes E$. The partial trace over $E$ is given
    entrywise by
    \begin{align}
        (\tr_E X)_{(a,c'),(a',d')}
        &=\sum_e X_{(a,c',e),(a',d',e)}.
        \notag
    \end{align}
    It preserves positivity and trace.
\end{lemma}

\begin{proof}
    \leanok
    Reindex $C$ as $C'\otimes E$ and apply the ordinary right partial trace.
    Its entry formula gives the displayed sum. Positivity follows by tracing
    a positive matrix, and invariance of the full trace under reindexing gives
    trace preservation.
\end{proof}

\begin{theorem}[Right marginals preserve a quantum Markov decomposition]
    \label{thm:quantum_markov_right_marginal}
    \lean{Entropy.rightMarginalAlong,
      Entropy.exists_quantumMarkovDecomposition_rightMarginalAlong}
    \leanok
    \uses{def:entropy_quantum_markov_decomposition,
      lem:partial_trace_right_along}
    Suppose that $C\simeq C'\otimes E$ and
    \begin{align}
        \rho_{ABC}
        &=\bigoplus_k p_k\rho_{A b_1}^{(k)}
          \otimes\rho_{b_2 C}^{(k)}
        \notag
    \end{align}
    is a quantum Markov decomposition on $B$. Then the marginal
    $\rho_{ABC'}=\tr_E(\rho_{ABC})$ has the decomposition
    \begin{align}
        \rho_{ABC'}
        &=\bigoplus_k p_k\rho_{A b_1}^{(k)}
          \otimes\tr_E\!\left(\rho_{b_2 C}^{(k)}\right).
        \notag
    \end{align}
\end{theorem}

\begin{proof}
    \leanok
    \uses{lem:partial_trace_right_along}
    The middle-system decomposition, its unitary change of basis, and the
    probabilities $p_k$ are unchanged. Partial trace preserves positivity and
    trace, so each $\tr_E(\rho_{b_2 C}^{(k)})$ is a density operator. Applying
    $\tr_E$ to the block-diagonal identity gives the displayed formula.
\end{proof}

\begin{lemma}[Tracing a terminal block]
    \label{lem:mpdo_reduced_state_collapse_last}
    \lean{collapse_last, collapse_last', reducedBlockState_cast}
    \leanok
    Let $a,b,c,N\in\mathbb N$ satisfy $a+b+c\leq N$. For configurations
    $u,v$ on the first $a+b$ sites,
    \begin{align}
        \sum_y \sigma_{a+b+c}^{(N)}(K)_
          {u\mathbin\Vert y,\,v\mathbin\Vert y}
        &=\sigma_{a+b}^{(N)}(K)_{u,v},
        \notag
    \end{align}
    where the sum ranges over configurations on the last $c$ sites. The same
    identity holds after replacing an arithmetically equal block length by its
    canonical identification.
\end{lemma}

\begin{proof}
    \leanok
    Expand both reduced states as partial traces of the normalized periodic
    state. The left-hand side first traces the complement of the
    $(a+b+c)$-site block and then its last $c$ sites; combining the two sums
    gives the partial trace onto the first $a+b$ sites.
\end{proof}

\begin{theorem}[SAL three-site marginals have Hayashi decompositions]
    \label{thm:mpdo_sal_three_site_eta_structure}
    \lean{MPOTensor.exists_etaStructure_reducedBlockState_three_of_isSAL}
    \leanok
    \uses{def:is_sal, def:mpdo_eta_structure}
    Let $K$ satisfy the strong area law and let $N\geq4$. If
    $\sigma_3^{(N)}(K)$ is obtained from the normalized $N$-site periodic state
    by tracing out all but its first three sites, then
    $\sigma_3^{(N)}(K)$ admits a quantum Markov decomposition on its middle
    site. This is the three-site marginal result proved in
    \cite[Appendix~C.2, lines~1351--1371]{Cirac2016MPDO_arXiv}. The
    decomposition may depend on $N$.
\end{theorem}

\begin{proof}
    \leanok
    \uses{thm:mpdo_sal_ssa_equality, thm:ssa_equality_reindex_factors,
        thm:sal_implies_eta_structure,
        thm:quantum_markov_right_marginal,
        lem:mpdo_reduced_state_collapse_last}
    Divide the first $N-1$ sites into consecutive regions of lengths $1$, $1$,
    and $N-3$. Let $\rho^{\mathrm{flat}}_{ABC}$ denote this marginal in the
    corresponding index set
    $\Fin{d^1}\times\Fin{d^1}\times\Fin{d^{N-3}}$.
    By Theorem~\ref{thm:mpdo_sal_ssa_equality},
    \begin{align}
        S(\rho^{\mathrm{flat}}_{ABC})+S(\rho^{\mathrm{flat}}_B)
        &=S(\rho^{\mathrm{flat}}_{AB})+S(\rho^{\mathrm{flat}}_{BC}).
        \notag
    \end{align}
    The first two singleton blocks have the canonical identification
    $e_{\mathrm{site}}:\Fin{d}\simeq\Fin{d^1}$, and reindexing these factors
    preserves the displayed equality. Positivity and unit trace follow from
    the normalized reduced state, so the Hayashi characterization gives a
    quantum Markov decomposition of the $(N-1)$-site marginal. Finally split
    the third region as $C\simeq c\otimes E$, where $c$ is its first site and
    $E$ contains the remaining $N-4$ sites.
    Theorem~\ref{thm:quantum_markov_right_marginal} gives the same middle-site
    direct sum after tracing out $E$. The contiguous-block reduction gives
    \begin{align}
        \sigma_3^{(N)}(K)
        &=\tr_{4,\ldots,N-1}\
          \left[\sigma_{N-1}^{(N)}(K)\right],
        \notag
    \end{align}
    which identifies the resulting operator with $\sigma_3^{(N)}(K)$.
\end{proof}

\begin{theorem}[Four-site SAL marginals have Hayashi decompositions]
    \label{thm:mpdo_four_site_sal_eta_structure}
    \lean{MPOTensor.exists_etaStructure_reducedBlockState_of_isSAL}
    \leanok
    \uses{def:is_sal, def:mpdo_eta_structure}
    Let $K$ satisfy the strong area law. If $\sigma_3^{(4)}(K)$ is obtained
    from its normalized four-site periodic state by tracing out the fourth
    site, then $\sigma_3^{(4)}(K)$ admits a quantum Markov decomposition on its
    middle site.
\end{theorem}

\begin{proof}
    \leanok
    \uses{thm:mpdo_sal_three_site_eta_structure}
    This is the case $N=4$ of
    Theorem~\ref{thm:mpdo_sal_three_site_eta_structure}.
\end{proof}

\begin{theorem}[BNT sector projectors for the four-site SAL marginal]
    \label{thm:mpdo_four_site_sal_bnt_sector_projectors}
    \lean{MPOTensor.reducedBlockState_four_reindex_eq_submatrix}
    \lean{MPOTensor.%
      exists_etaStructure_reducedBlockState_four_reindex_of_isSAL}
    \lean{MPOTensor.exists_bntSectorProjectors_four_of_sameMPV₂Pos_isSAL}
    \leanok
    \uses{thm:mpdo_normalized_three_site_bnt_family_closure,
      thm:mpdo_bnt_sector_projection_resolution,
      thm:mpdo_four_site_sal_eta_structure}
    Suppose that the common blocking has been performed, so that the BNT
    representatives have a simultaneous one-letter span. For the normalized
    three-site marginal $\sigma_3^{(4)}$, there are a simultaneous inverse
    $C$ and a Hayashi decomposition such that every Markov sector $k$ with
    $p_k\ne0$ has a unique BNT label $j_k$. Consequently,
    \begin{align}
        P_s
        &=\sum_{\substack{k:p_k\ne0\\j_k=s}}Q_k.
        \notag
    \end{align}
    Each $P_s$ is an orthogonal projection,
    $P_sP_t=0$ for $s\ne t$, and $\sum_sP_s=P_{\mathrm{act}}$.
    No independent nonzero-closing-matrix hypothesis is required. This is
    the four-site specialization of equations~\textup{QkKjs} and
    \textup{Pis} in
    \cite[Appendix~C.2, lines~1714--1737]{Cirac2016MPDO_arXiv}.
    The Hayashi decomposition uses the forward equality characterization of
    strong subadditivity recorded in
    Theorem~\ref{thm:hayashi_ssa_equality_characterization}; no additional
    analytic assumption is introduced here.
\end{theorem}

\begin{proof}\leanok
    The canonical ordered-triple reindexing is the submatrix convention of
    the SAL-to-Hayashi theorem. Choose the simultaneous inverse from the
    one-letter span and the Hayashi decomposition from SAL. By
    Theorem~\ref{thm:mpdo_normalized_three_site_bnt_family_closure}, the same
    marginal is a BNT family closure whose closing matrices are all nonzero.
    The generic uniqueness and projector theorems now give the unique labels,
    orthogonality, mutual disjointness, and resolution of the active support.
\end{proof}

\begin{definition}[Sitewise physical action and absorbed BNT representatives]
    \label{def:mpdo_sitewise_physical_matrix}
    \lean{MPOTensor.sitewisePhysicalMatrix}
    \lean{MPOTensor.commonWeightAbsorbedBasisMPOTensor}
    \leanok
    For a one-site matrix $V$, write $V^{\otimes N}$ in the configuration
    basis as
    \begin{align}
        (V^{\otimes N})_{s,t}
        &=\prod_{n=1}^{N}V_{s_n,t_n}.
        \notag
    \end{align}
    If the copies of the $j$-th normal tensor have common coefficient
    $\mu_j$, let $\mathcal K_j=\mu_j A_j$ denote the representative with this
    coefficient absorbed into its local tensor.
\end{definition}

\begin{theorem}[Local selection by the BNT projectors]
    \label{thm:mpdo_bnt_projector_local_selection}
    \lean{MPOTensor.bntSectorProjection_mul_physicalSlice_mul_eq_zero}
    \lean{MPOTensor.bntSectorProjection_mul_physicalSlice_mul_self}
    \lean{MPOTensor.%
      sum_bntSectorProjection_mul_physicalSlice_mul_self}
    \lean{MPOTensor.changePhysicalBasis_bntSectorProjection_basis}
    \lean{MPOTensor.%
      exists_bntProjectorSelection_positiveLength_of_sameMPV₂Pos_isSAL}
    \leanok
    \uses{thm:mpdo_four_site_sal_bnt_sector_projectors}
    Under the hypotheses of
    Theorem~\ref{thm:mpdo_four_site_sal_bnt_sector_projectors}, the resulting
    projectors satisfy $P_sA_iP_t=0$ whenever $s\ne t$ or $i\ne s$.
    For the matching label, they satisfy $P_iA_iP_i=A_i$. Consequently,
    $\sum_sP_sA_iP_s=A_i$, which is the local diagonal decomposition required
    for the sector entropy argument.
    Equivalently, compressing the physical legs of $A_i$ by $P_s$ retains
    $A_i$ for $s=i$ and gives the zero local tensor otherwise. This is the
    local projection relation in
    \cite[Appendix~C.2, lines~1680--1755]{Cirac2016MPDO_arXiv}.

    The Markov summands of zero Hayashi weight are retained in the ambient
    decomposition. They annihilate every normal-sector physical slice, so
    the displayed local identities do not require the stronger assertion
    $\sum_sP_s=\Id$ outside the active support.
\end{theorem}

\begin{proof}\leanok
    Distinct Hayashi sectors give zero off-diagonal corners. On an active
    diagonal sector, uniqueness of the BNT label kills every normal tensor
    except the matching one. A zero-weight diagonal sector also kills every
    physical slice: otherwise the nonzero closing matrix would make the
    corresponding three-site Markov block nonzero. Summing these statements
    over the Hayashi sectors gives the asserted $P_sA_iP_t$ identities.
\end{proof}

\begin{definition}[Completion on the zero-weight Markov sectors]
    \label{def:mpdo_completed_bnt_sector_projection}
    \lean{MPOTensor.completedBntSectorLabel}
    \lean{MPOTensor.completedBntSectorProjection}
    \leanok
    \uses{def:mpdo_bnt_sector_projection}
    Fix a BNT label $s_0$. Each positive-weight Markov sector retains its
    unique BNT label, while every zero-weight Markov sector is assigned the
    label $s_0$. If $\widehat S_s$ is the set of Markov sectors with completed
    label $s$, define
    \begin{align}
        \widehat P_s
        &=\sum_{k\in\widehat S_s}Q_k.
        \notag
    \end{align}
\end{definition}

\begin{theorem}[Separating projectors with copy-independent weights]
    \label{thm:mpdo_sal_bnt_separating_projectors_copy_independent}
    \lean{MPOTensor.completedBntSectorProjection_isOrthogonal}
    \lean{MPOTensor.completedBntSectorProjection_mul_eq_zero}
    \lean{MPOTensor.sum_completedBntSectorProjection}
    \lean{MPOTensor.%
      completedBntSectorProjection_mul_physicalSlice_mul_eq_zero}
    \lean{MPOTensor.%
      exists_bntSeparatingProjectors_of_sameMPV₂Pos_isSAL_of_weight_copy_independent}
    \leanok
    \uses{def:mpdo_completed_bnt_sector_projection,
      thm:mpdo_bnt_projector_local_selection}
    Let $K$ be simple, in biCF, and satisfy SAL. Assume in addition that the
    common blocking has been performed, the copy weights over each BNT
    representative agree, and the representatives have a simultaneous
    one-site span. There is a distinguished BNT label $s_0$ such that the
    completed projectors are orthogonal and, for distinct labels $s$ and $t$,
    \begin{align}
        \widehat P_s\widehat P_t&=0,
        \notag\\
        \sum_s\widehat P_s&=\Id.
        \notag
    \end{align}
    For every physical slice of the $i$-th BNT representative, whenever
    $s\ne t$ or $i\ne s$,
    \begin{align}
        \widehat P_sK_i^{\beta\alpha}\widehat P_t&=0.
        \notag
    \end{align}
    This is the copy-independent-weight specialization of the
    separating-projector lemma in
    \cite[Appendix~C.2, lines~1680--1737]{Cirac2016MPDO_arXiv}.
\end{theorem}

\begin{proof}
    \leanok
    \uses{thm:mpdo_four_site_sal_bnt_sector_projectors,
      thm:mpdo_bnt_projector_local_selection}
    A unit-modulus BNT weight supplies the label $s_0$. The Markov-sector
    projections $Q_k$ are mutually orthogonal and sum to the identity, so
    partitioning all labels by $\widehat S_s$ proves orthogonality and the
    identity resolution. If $p_k=0$, then
    $Q_kK_i^{\beta\alpha}Q_k=0$ for every $i,\beta,\alpha$. If $p_k\ne0$, then
    $Q_kK_i^{\beta\alpha}Q_k=0$ whenever the BNT label of sector $k$ is not
    $i$. For distinct Markov sectors,
    $Q_kK_i^{\beta\alpha}Q_l=0$ whenever $k\ne l$. Hence, whenever $s\ne t$
    or $i\ne s$,
    \begin{align}
        \widehat P_sK_i^{\beta\alpha}\widehat P_t
        &=\sum_{k\in\widehat S_s}\sum_{l\in\widehat S_t}
          Q_kK_i^{\beta\alpha}Q_l
          =0,
        \notag
    \end{align}
    since every summand vanishes by one of the three preceding cases.
\end{proof}

\begin{theorem}[Separating projectors for a simple tensor satisfying ZCL and SAL]
    \label{thm:mpdo_sal_bnt_separating_projectors}
    \lean{MPOTensor.%
      exists_bntSeparatingProjectors_of_horizontalCF_isSourceZCL_isSAL}
    \leanok
    \uses{thm:mpdo_simple_weight_copy_independent,
      thm:mpdo_sal_bnt_separating_projectors_copy_independent}
    Let $K$ be a simple tensor in block-injective canonical form which satisfies
    ZCL and SAL. For the representatives $\mathcal K_i$ of its BNT, there are
    orthogonal projectors $P_s$ such that
    \begin{align}
        \sum_sP_s&=\Id,
        \notag\\
        P_sK_i^{\beta\alpha}P_t&=0
        \quad\text{whenever }s\ne t\text{ or }i\ne s.
        \notag
    \end{align}
    This is the separating-projector lemma in
    \cite[Appendix~C.2, lines~1626--1691]{Cirac2016MPDO_arXiv}.
\end{theorem}

\begin{proof}
    \leanok
    Write the horizontal canonical form as
    $K=G(\bigoplus_{j,q}\mu_{j,q}\mathcal K_j)G^{-1}$. Restricting the ZCL
    equation to the $(j,q)$-th block and using the non-nilpotence supplied by
    simplicity gives $\mu_{j,q}=\mu_{j,q'}$ for all $j,q,q'$.
    The block-diagonal gauge preserves every positive-length periodic state,
    so the canonical-form tensor and $K$ generate the same such states. The
    preceding copy-independent-weight theorem now supplies the projectors.
\end{proof}

\begin{theorem}[Positive-length BNT-sector compression]
    \label{thm:mpdo_bnt_projector_positive_length_compression}
    \lean{MPOTensor.singleKrausMap_sitewise_bntSectorProjection_mpo}
    \lean{MPOTensor.sitewise_bntSectorProjection_mul_mpo_mul_self}
    \lean{MPOTensor.%
      exists_bntProjectorSelection_positiveLength_of_sameMPV₂Pos_isSAL}
    \leanok
    \uses{def:mpdo_sitewise_physical_matrix,
      thm:mpdo_bnt_projector_local_selection}
    For every nonempty periodic chain and every normal-sector label $s$,
    \begin{align}
        P_s^{\otimes N} \sigma^{(N)}(K)\,P_s^{\otimes N}
        &=n_s\,\sigma^{(N)}(\mathcal K_s),
        \qquad N\geq1.
        \label{eq:rfp_positive_compression}
    \end{align}
    This is equation~\textup{sigmaNKj} in
    \cite[Appendix~C.2, lines~1756--1759]{Cirac2016MPDO_arXiv}.

    Here and throughout the MPDO discussion, a periodic chain has positive
    length.
\end{theorem}

\begin{proof}\leanok
    Sitewise multiplication by $P_s$ is the periodic MPO obtained by changing
    every local physical matrix by $P_s$. The preceding theorem replaces
    each normal representative by itself when its label is $s$ and by the
    zero tensor otherwise. For $N\geq1$, the latter has zero periodic MPO.
    The surviving coefficient is
    $n_s\mu_s^N$, which is precisely $n_s$ times the MPO of the representative
    with $\mu_s$ absorbed locally.
\end{proof}

\begin{corollary}[Each absorbed BNT element generates MPDOs]
    \label{cor:mpdo_absorbed_bnt_sector_is_mpdo}
    \lean{MPOTensor.%
      commonWeightAbsorbedBasisMPOTensor_isMPDO_of_projectorSelection}
    \lean{MPOTensor.%
      commonWeightAbsorbedBasisMPOTensor_isMPDO_of_sameMPV₂Pos_isSAL}
    \leanok
    \uses{thm:mpdo_bnt_projector_positive_length_compression}
    Under the hypotheses of
    Theorem~\ref{thm:mpdo_bnt_projector_positive_length_compression}, every
    absorbed BNT element $\mathcal K_s$ generates positive semidefinite
    operators on all nonempty periodic chains. Indeed,
    $P_s^{\otimes N}\sigma^{(N)}(K)P_s^{\otimes N}\geq0$ and $n_s>0$, so
    (\ref{eq:rfp_positive_compression}) gives
    $\sigma^{(N)}(\mathcal K_s)\geq0$. This is the first positivity conclusion
    in~\cite[Appendix~C.2, lines~1753--1759]{Cirac2016MPDO_arXiv}.
\end{corollary}

\begin{proof}\leanok
    Congruence by the Hermitian projection $P_s^{\otimes N}$ preserves
    positive semidefiniteness. Multiplication by the inverse of the positive
    integer $n_s$ then gives positivity of the sector operator.
\end{proof}

% Source: arXiv:1606.00608, Appendix C.2, lines 1745--1782.
\begin{theorem}[Literal ZCL descends to absorbed BNT representatives]
    \label{thm:mpdo_absorbed_bnt_sector_literal_zcl}
    \lean{MPOTensor.weighted_basis_physTraceTransfer_sq_of_literal_ZCL}
    \lean{MPOTensor.%
      physTraceTransfer_commonWeightAbsorbedBasisMPOTensor}
    \lean{MPOTensor.%
      commonWeightAbsorbedBasisMPOTensor_physTraceTransfer_sq_of_literal_ZCL}
    \leanok
    \uses{def:mpo_physical_trace_transfer,
      def:mpdo_common_sector_weight_absorption}
    Let $K$ be gauge equivalent to the weighted direct sum of BNT
    representatives $A_s$, with copy weights $\mu_{s,q}$. If its
    physical-trace transfer is literally idempotent, then every weighted copy
    satisfies
    \begin{align}
        \mathcal T_{\mu_{s,q}A_s}^{2}
        &=\mathcal T_{\mu_{s,q}A_s}.
        \notag
    \end{align}
    If the copy weights over $A_s$ agree and
    $\mathcal K_s=\mu_sA_s$, then
    \begin{align}
        \mathcal T_{\mathcal K_s}^{2}
        &=\mathcal T_{\mathcal K_s}.
        \notag
    \end{align}
\end{theorem}

\begin{proof}\leanok
    \uses{def:global_gauge_of_blocks,
      lem:mpdo_doubled_phys_trace_transfer_to_mps,
      lem:mpdo_doubled_phys_trace_transfer_blockwise}
    Conjugating $\mathcal T_K^2=\mathcal T_K$ by the block-diagonal gauge
    gives the same identity for the weighted direct sum. Its diagonal block
    indexed by $(s,q)$ is
    $\mathcal T_{\mu_{s,q}A_s}^{2}
      =\mathcal T_{\mu_{s,q}A_s}$.
    When the weights are copy-independent,
    $\mathcal T_{\mathcal K_s}=\mu_s\mathcal T_{A_s}$ is that block for any
    copy over $s$.
\end{proof}

% Source: arXiv:1606.00608, Appendix C.2, Lemma `lemmus`, lines
% 1646--1661, and the absorbed-sector ZCL conclusion at line 1781.
\begin{theorem}[Scale-invariant source ZCL descends to absorbed BNT representatives]
    \label{thm:mpdo_absorbed_bnt_sector_zcl}
    \lean{MPOTensor.IsSimpleCanonicalForm.%
      exists_commonWeightAbsorbedBasisMPOTensor_isSourceZCL}
    \lean{MPOTensor.commonWeightAbsorbedBasisMPOTensor_isSourceZCL}
    \leanok
    \uses{def:mpdo_simple_cf_tensor, def:mpo_source_zcl,
      def:mpdo_common_sector_weight_absorption}
    Let $K$ be a simple tensor in its chosen canonical form and suppose that
    $K$ has scale-invariant source zero correlation length. Then the copy
    weights over each normal representative $A_s$ agree. For every $A_s$ with
    non-nilpotent physical-trace transfer $\mathcal B_s$, absorb the common
    weight and write $\mathcal K_s=\mu_sA_s$. Each absorbed BNT representative
    $\mathcal K_s$ has scale-invariant source zero correlation length. More
    precisely,
    \begin{align}
        \mathcal T_{\mathcal K_s}
        &=\mu_s\mathcal B_s\ne0,
        \notag\\
        \mathcal T_{\mathcal K_s}^{2}
        &=\lambda\mathcal T_{\mathcal K_s},
        \qquad \lambda>0.
        \notag
    \end{align}
    This is the final zero-correlation-length conclusion in
    \cite[Appendix~C.2, line~1781]{Cirac2016MPDO_arXiv}.
\end{theorem}

\begin{proof}\leanok
    \uses{cor:mpdo_simple_zcl_weight_copy_independent,
      lem:mpdo_simple_blockwise_zcl_equation}
    Corollary~\ref{cor:mpdo_simple_zcl_weight_copy_independent} gives the
    common weight on every copy. Lemma~\ref{lem:mpdo_simple_blockwise_zcl_equation}
    then gives the displayed quadratic identity. The common weight is nonzero,
    and nonnilpotence of $\mathcal B_s$ implies $\mathcal B_s\ne0$, so the
    absorbed transfer is nonzero.
\end{proof}

\begin{remark}[Sectorwise hypotheses in the canonical BNT application]
    \label{rem:mpdo_prop4to2_sectorwise_zcl_boundary}
    \uses{thm:mpdo_absorbed_bnt_sector_zcl,
      thm:mpdo_bnt_physical_support_restriction,
      thm:mpdo_proportional_commuting_sector_positivity,
      thm:mpdo_canonical_bnt_proportional_sectors_zcl_sal}
    In the canonical absorbed BNT presentation, ambient source zero
    correlation length supplies source zero correlation length for every
    absorbed representative by
    Theorem~\ref{thm:mpdo_absorbed_bnt_sector_zcl}. The standing biCF
    one-letter span and the nonzero common weights give one-site injectivity by
    Theorem~\ref{thm:mpdo_bnt_physical_support_restriction}.

    The positive commuting-bond realization begins at length two. It therefore
    proves positivity of $\rho^{(N)}(K_x)$ for all $N\geq2$, while positivity
    of $\rho^{(1)}(K_x)$ remains the only separate sector-positivity
    hypothesis. Theorem~\ref{thm:mpdo_canonical_bnt_proportional_sectors_zcl_sal}
    combines these facts and proves SAL at every original cut, without a
    normality hypothesis on the absorbed sectors. The printed GSNNCH-only
    proposition remains incomplete because it omits ambient source zero
    correlation length; an arbitrary GSNNCH presentation need not inherit this
    tensor-level property.

    % **Scope restriction (ambient source ZCL and one-site positivity):**
    % Proposition `prop4to2` omits ambient source ZCL, and its commuting-bond
    % identity begins at length two. See
    % docs/paper-gaps/cpsv16_commuting_form_to_sal.tex.
\end{remark}

\begin{theorem}[Strict normalization at positive length]
    \label{thm:mpdo_source_zcl_strict_positive_length_normalization}
    \lean{MPOTensor.trace_mpo_ne_zero_of_isSourceZCL}
    \lean{MPOTensor.trace_mpo_pos_of_isMPDO_isSourceZCL}
    \lean{MPOTensor.%
      trace_mpo_commonWeightAbsorbedBasisMPOTensor_pos}
    \leanok
    \uses{def:mpdo, def:mpo_source_zcl, lem:mpdo_trace_loop_chain,
      cor:mpdo_absorbed_bnt_sector_is_mpdo,
      thm:mpdo_absorbed_bnt_sector_zcl}
    Let $M$ generate matrix product density operators and have source zero
    correlation length. Then, for every $N\geq1$,
    $\tr(\sigma^{(N)}(M))>0$.
    In particular, after a common blocking has made the simultaneous BNT
    word span full at one letter, the absorbed BNT elements selected in
    Theorem~\ref{thm:mpdo_absorbed_bnt_sector_zcl} have a well-defined
    normalized state at every physical chain length. This is the biCF
    hypothesis imposed at the start of Case II in the source; its relation
    with finite physical blocking is recorded in
    \path{docs/paper-gaps/cpgsv17_bicf_block_separation.tex}.
    This supplies the strict form of the sector normalization used in
    \cite[Appendix~C.2, lines~1760--1780]{Cirac2016MPDO_arXiv}.
\end{theorem}

\begin{proof}\leanok
    Write the physical-trace transfer as $\mathcal T_M=\lambda E$, where
    $\lambda>0$ and $E$ is a nonzero idempotent. A nonzero idempotent on a
    finite-dimensional complex vector space has nonzero trace, while
    $E^N=E$ for $N\geq1$. Hence
    \begin{align}
        \tr(\sigma^{(N)}(M))
        &=\tr(\mathcal T_M^N)
          =\lambda^N\tr(E)
          \ne0.
        \notag
    \end{align}
    Positivity of the density operator then makes this trace strictly
    positive.
\end{proof}

\begin{theorem}[Saturated area law for the absorbed BNT elements]
    \label{thm:mpdo_absorbed_bnt_sector_sal}
    \lean{MPOTensor.bntSectorProbability}
    \lean{MPOTensor.bntSectorProbability_pos}
    \lean{MPOTensor.normalizedMPO_eq_sum_bntSectorProbability_smul}
    \lean{MPOTensor.reducedBlockState_eq_sum_bntSectorProbability_smul}
    \lean{MPOTensor.blockEntropy_eq_sum_bntSectorProbability}
    \lean{MPOTensor.blockEntropy_add_blockEntropy_eq_bntSector}
    \lean{MPOTensor.%
      commonWeightAbsorbedBasisMPOTensor_isSAL_of_projectorSelection}
    \lean{MPOTensor.%
      commonWeightAbsorbedBasisMPOTensor_isSAL_of_sameMPV₂Pos}
    \leanok
    \uses{def:is_sal, def:mpo_source_zcl,
      def:mpdo_common_sector_weight_absorption}
    Let $K$ saturate the area law and have zero correlation length, and let
    its horizontal BNT canonical form, up to block-diagonal gauge, have normal
    representatives $A_s$, common copy weights $\mu_s$, and copy numbers
    $n_s$. Assume that the physical-trace transfer of every $A_s$ is
    nonnilpotent. After a common blocking, assume explicitly that the
    simultaneous one-letter word tuples of the $A_s$ span the direct sum of
    their full matrix algebras. Put $\mathcal K_s=\mu_sA_s$.

    For every nonempty chain, define
    \begin{align}
        p_s^{(N)}
        &=
        \frac{n_s\,\tr\!\left(\sigma^{(N)}(\mathcal K_s)\right)}
             {\tr\!\left(\sigma^{(N)}(K)\right)}.
        \notag
    \end{align}
    Then $p_s^{(N)}>0$, and every nonempty prefix marginal has the same sector
    decomposition with these probabilities. Translation invariance
    identifies the other contiguous blocks of the same length.
    Consequently, every absorbed representative $\mathcal K_s$ saturates the
    area law.

    The one-letter simultaneous span is the block-injective canonical-form
    hypothesis imposed at the start of Case II in
    \cite[line~1628]{Cirac2016MPDO_arXiv}. Its relation with finite physical
    blocking is recorded in
    \path{docs/paper-gaps/cpgsv17_bicf_block_separation.tex}. The entropy
    conclusion is the argument of
    \cite[Appendix~C.2, lines~1748--1781]{Cirac2016MPDO_arXiv}.
\end{theorem}

\begin{proof}\leanok
    \uses{thm:entropy_pairwise_annihilating_support,
      thm:positive_weighted_sum_termwise_equality,
      lem:sal_const, thm:entropy_strong_subadditivity,
      thm:mpdo_tripartite_block_entropy_identifications,
      thm:mpdo_bnt_projector_local_selection,
      cor:mpdo_absorbed_bnt_sector_is_mpdo,
      thm:mpdo_source_zcl_strict_positive_length_normalization}
    The BNT projections carry every normalized marginal on mutually
    annihilating supports, so the preceding support form of entropy
    additivity gives
    \begin{align}
        S_L(K)
        &=H\!\left(p^{(N)}\right)
          +\sum_s p_s^{(N)}S_L(\mathcal K_s),
        \qquad 1\leq L\leq N.
        \notag
    \end{align}
    The four marginals entering strong subadditivity have the same
    probabilities. Substitution in the equality for $K$ cancels all four
    Shannon terms. Strong subadditivity gives the corresponding inequality
    in each sector, and strict positivity of every $p_s^{(N)}$ forces every
    such inequality to be an equality. Comparing the equalities for block
    lengths $L$ and $L+1$ proves saturation for $\mathcal K_s$.
\end{proof}

% Source: arXiv:1606.00608, Appendix C.2, lines 1740--1782.
\begin{theorem}[Analytic properties of absorbed BNT representatives]
    \label{thm:mpdo_absorbed_bnt_sector_caseii_analytic_properties}
    \lean{MPOTensor.%
      commonWeightAbsorbedBasisMPOTensor_caseII_properties_of_literal_ZCL}
    \leanok
    \uses{def:injective, def:mpdo, def:is_sal,
      def:mpo_physical_trace_transfer,
      def:mpdo_common_sector_weight_absorption}
    Let $K$ be an MPDO tensor in biCF, written up to a block-diagonal
    gauge as a weighted BNT direct sum whose BNT representatives have
    nonnilpotent physical-trace transfers. Suppose that the weights over each
    BNT label are copy-independent,
    $\mu_{s,q}=\mu_{s,q'}$, and that $K$ saturates the area law and satisfies
    literal zero correlation length, $\mathcal T_K^2=\mathcal T_K$. For every
    BNT label $s$, let $\mathcal K_s=\mu_sA_s$ be the absorbed BNT
    representative. Then
    $\mathcal K_s$ is one-site injective, generates positive semidefinite
    operators on every nonempty periodic chain, saturates the area law, and
    satisfies
    \begin{align}
        \mathcal T_{\mathcal K_s}^{2}
        &=\mathcal T_{\mathcal K_s}.
        \notag
    \end{align}
\end{theorem}

\begin{proof}\leanok
    \uses{def:mpo_physical_trace_transfer,
      lem:mpdo_trace_loop_chain,
      thm:mpo_source_zcl_of_physical_trace_idempotent,
      thm:mpdo_bnt_physical_support_restriction,
      cor:mpdo_absorbed_bnt_sector_is_mpdo,
      thm:mpdo_absorbed_bnt_sector_sal,
      thm:mpdo_absorbed_bnt_sector_literal_zcl}
    Saturation gives $\tr(\sigma^{(1)}(K))>0$. Since
    $\tr(\sigma^{(1)}(K))=\tr(\mathcal T_K)$, one has
    $\mathcal T_K\ne0$. Together with $\mathcal T_K^2=\mathcal T_K$, this
    gives scale-invariant source zero correlation length with scalar $1$.
    The one-letter BNT span gives injectivity, the sector compression gives
    positivity, and the sector entropy decomposition gives saturation of the
    area law. Theorem~\ref{thm:mpdo_absorbed_bnt_sector_literal_zcl} supplies
    the final literal idempotence identity.
\end{proof}
